\section{The genuine graded quotient after strict parameter elimination}

Let $R$ be a commutative ring, $I,N$ ideals, and $v\in N$. Write
$D_q(I)=\operatorname{im}(N^q\to R/(I+N^{q+1}))$. These are the actual
power-image layers, with their inherited $R$-module structures. Throughout
$I$ and $N$ are unchanged entire ideals. Suppose, for every $q\geq0$,
$b\in N^q$ and $vb\in I+N^{q+2}$ imply $b\in I+N^{q+1}$.

\begin{theorem}[The canonical layer map is onto]
\label{hir:strict-parameter-layer-map-surjective}
\lean{Hironaka.HironakaStrictFilteredGradedQuotient.adicLayer,Hironaka.HironakaStrictFilteredGradedQuotient.adicLayerMap,Hironaka.HironakaStrictFilteredGradedQuotient.adicLayerMap_surjective}
\leanok
The original quotient factor defines a canonical surjective map
$D_q(I)\to D_q(I+(v))$ in every degree. It sends each original numerator
$b\in N^q$ to the class of that same $b$.
\end{theorem}
\begin{proof}
The inclusion $I+N^{q+1}\subset I+(v)+N^{q+1}$ defines the quotient
factor. Its image on a class represented by $b\in N^q$ lies in the
claimed target image. Every target layer element has just such an
original representative $b\in N^q$, and its old class is a preimage.
Both maps use the original quotient scalar actions.
\end{proof}

\begin{theorem}[Actual multiplication gives the degree-one symbol map]
\label{hir:strict-parameter-symbol-map}
\lean{Hironaka.HironakaStrictFilteredGradedQuotient.adicSymbolMul,Hironaka.HironakaStrictFilteredGradedQuotient.adicSymbolMul_apply}
\leanok
Multiplication by the original element $v$ gives an $R$-linear map
$D_q(I)\to D_{q+1}(I)$, sending each original numerator $b$ to $vb$.
\end{theorem}
\begin{proof}
If $b-b'\in I+N^{q+1}$, multiplying by $v$ gives
$v(b-b')\in I+N^{q+2}$, because $vI\subset I$ and
$vN^{q+1}\subset N^{q+2}$. Thus multiplication descends from the entire
quotient. For $b\in N^q$, its image $vb$ lies in $N^{q+1}$, so this
descended map restricts to the claimed power-image layers. Multiplication
by a fixed element of a commutative ring is $R$-linear.
\end{proof}

\begin{theorem}[The entire positive-degree kernel is the actual symbol image]
\label{hir:strict-parameter-whole-layer-kernel}
\uses{hir:strict-parameter-layer-map-surjective,hir:strict-parameter-symbol-map,hir:regular-symbol-whole-principal-intersection}
\lean{Hironaka.HironakaStrictFilteredGradedQuotient.ker_adicLayerMap_eq_range_adicSymbolMul,Hironaka.HironakaStrictFilteredGradedQuotient.adicLayerQuotientEquiv,Hironaka.HironakaStrictFilteredGradedQuotient.adicLayerQuotientEquiv_apply}
\leanok
In every degree $q+1$, the kernel of the canonical map to
$D_{q+1}(I+(v))$ is exactly the image of
$v:D_q(I)\to D_{q+1}(I)$. Consequently the quotient by this entire image
is canonically $R$-linearly isomorphic to $D_{q+1}(I+(v))$.
\end{theorem}
\begin{proof}
Take a kernel element represented by $b\in N^{q+1}$. Its vanishing in
the target means $b=i+va+c$ with $i\in I$ and $c\in N^{q+2}$.
Thus $b-c$ belongs to both $I+N^{q+1}$ and $I+(v)$. The strict
intersection theorem gives $b-c=j+vh$ with $j\in I$ and $h\in N^q$.
Modulo $I+N^{q+2}$ the entire old class of $b$ is therefore the image
of the old degree-$q$ class of $h$ under actual multiplication by $v$.
Conversely each such image is killed by the quotient imposing $v$.
This proves equality of the whole kernel and image. The first
isomorphism theorem and the proved surjectivity give the stated
quotient equivalence, using the actual map and unchanged scalar action.
No separatedness is required.
\end{proof}

\begin{theorem}[Degree zero is unchanged]
\label{hir:strict-parameter-degree-zero}
\lean{Hironaka.HironakaStrictFilteredGradedQuotient.adicLayer_zero_denominator,Hironaka.HironakaStrictFilteredGradedQuotient.adicLayerMap_zero_injective}
\leanok
The degree-zero denominators agree: $I+(v)+N=I+N$.
\end{theorem}
\begin{proof}
Since $v\in N$, the entire principal ideal $(v)$ is contained in $N$.
Both inclusions now follow by associativity of ideal sums. If two degree-zero source classes have equal images, their original
numerator difference belongs to the target denominator, hence to the
same old denominator by the displayed equality. The canonical map is
therefore injective and, by the first theorem, surjective. Thus degree
zero agrees under the canonical quotient factor, as required alongside
the positive-degree quotient equivalences.
\end{proof}
